\chapter{Post-processing}
\label{ch:postprocessing}
\index{post-processing}

\section{Output families}
GEOS-MPM cases commonly produce plot files, restart files, PFW logs, XML input, particle files, scheduler scripts, and optional CSV histories.  Plot output is selected with the PFW \texttt{outputType} key or directly in XML.

\section{VisIt}
\index{VisIt}
Silo output is the standard VisIt pathway.  Open the \path{output_*} file set and add mesh plots for the grid or particles, or pseudocolor/vector plots for particle and requested grid variables.  The suite workflow includes \path{pfw_visit_render.py}, a command-line renderer for initial/final frames after a case has produced Silo or VTK output.

\begin{lstlisting}[language=bash,caption={VisIt batch rendering helper from the suite workflow.}]
visit -nowin -cli -s pfw_visit_render.py \
  --run-dir /p/lustre1/$USER/geos-mpm-suite-runs/verification/Ftable/elasticBlockUni \
  --variable Damage \
  --view xy
\end{lstlisting}

\section{ParaView}
\index{ParaView}
VTK output is the ParaView pathway.  Configure and build GEOS with VTK output enabled, then set \texttt{pfw["outputType"] = "vtk"}.  ParaView can load the generated VTK file sequence for particle and mesh variables.  Use VTK output for ParaView-oriented exploration; use Silo output for workflows that require the current VisIt scripts or Silo-native grouped vector variables.

\section{Analysis scripts}
\index{analysis scripts}
The PFW tree contains suite-level and case-specific analysis scripts.  They include suite post-processing/reporting, rendered-frame generation, verification/validation plotting, reaction-history plots, material-model checks, and example-specific post-processors.  Appendix~\ref{app:post-scripts} lists the discovered Python and shell scripts in the current archive.

\section{CSV histories}
\label{sec:csv-histories}
\index{CSV histories}
Reaction, box-average, profile, and tracer histories are intentionally simple CSV outputs.  They are useful for quick Python, MATLAB, Julia, or shell-based checks and for suite status scans.  Because file names and headers are controlled by the solver, downstream scripts should validate header names rather than rely only on column positions.

The reaction file \path{reactionHistory.csv} has the header \texttt{time, F00, F11, F22, length\_x, length\_y, length\_z, Rx-, Rx+, Ry-, Ry+, Rz-, Rz+, L00, L11, L22}.  The reaction columns are integrated face forces in the solver force units.  To compare to a measured or analytical stress, divide the appropriate reaction by the actual material area of the loaded boundary.  The PFW controls that enable this file and select the velocity-correction or internal-force reaction measure are described in Section~\ref{subsec:pfw-reaction-history}.

The box-average file \path{boxAverageHistory.csv} reports volume- or mass-weighted aggregate quantities over the selected averaging box.  Profile files use the pattern \path{profile_<direction>_<variable>.csv}.  Tracer files use the configured \texttt{tracerOutputPrefix} and write \texttt{t,x,y,z} followed by the requested particle fields for the selected particle ID.

